\documentclass[UTF8,12pt,a4paper]{ctexart}

\usepackage{amsmath}
\usepackage{cases}
\usepackage{cite}
\usepackage{graphicx}
\usepackage{enumerate}
\usepackage{algorithm}
\usepackage{caption}   % \caption*需要
\usepackage{subcaption} % 子图布局
\usepackage[noend]{algpseudocode} %algorithmicx
\makeatletter
\renewcommand{\fnum@algorithm}{\fname@algorithm}
\makeatother
\renewcommand{\algorithmicrequire}{\textbf{Input:}}
\renewcommand{\algorithmicensure}{\textbf{Output:}}
\usepackage[margin=1in]{geometry}
\geometry{a4paper}
\usepackage{fancyhdr}
\pagestyle{fancy}
\fancyhf{}
\usepackage{xcolor}

% ------------------------- 代码展示设置 -------------------------
\usepackage{listings}
\lstset{
	basicstyle=\ttfamily,                                % 传统上，代码打字机字体好些
	breaklines,
	tabsize=4,
	extendedchars=false,
	columns=fixed,
	%numbers=left,                                        % 在左侧显示行号
	numbers=none,
	frame=none,                                          % 不显示背景边框
	backgroundcolor=\color[RGB]{245,245,244},            % 设定背景颜色
	keywordstyle=\color[RGB]{40,40,255},                 % 设定关键字颜色
	numberstyle=\footnotesize\color{darkgray},           % 设定行号格式
	commentstyle=\it\color[RGB]{0,96,96},                % 设置代码注释的格式
	stringstyle=\rmfamily\slshape\color[RGB]{128,0,0},   % 设置字符串格式
	showstringspaces=false,                              % 不显示字符串中的空格
	xleftmargin=1em,
	%xrightmargin=1em,
	%aboveskip=-0.5em
	language=c++,                                        % 设置语言
}



% ------------------- 设置超链接，便于跳转 -----------------
\usepackage{enumitem}
\usepackage{hyperref}
\hypersetup{
    pdfborder={0 0 0},    % 消除超链接边框
    colorlinks=true,       % 将边框改为颜色标记（可选）
    linkcolor=black,       % 内部链接颜色（目录、引用等）
    citecolor=black,       % 引用颜色
    urlcolor=blue          % URL链接颜色（保持蓝色以便区分）
}




% ------------------------ 标题区 ----------------------------
\title{算法原理2025-2026 期末试题}
\pagenumbering{arabic}

\begin{document}



% ------------------------ 页眉和页脚 ----------------------------

\fancyhead[C]{算法原理2025~2026 期末考试}
\fancyfoot[C]{\thepage}

\maketitle
%\tableofcontents
%\newpage




% ============================================================
% 第1题：概念题
% ============================================================
\section{试题}

% 使用 enumerate 环境，设置 itemsep 为一行高度
\begin{enumerate}[label=\textbf{\arabic*.}, itemsep=\baselineskip]

    \item 解释如下概念；
    \begin{enumerate}
      \item 回顾课程中介绍的递归方程求解的三种方法，每种方法进行简要说明
      \item 给出一个学习生活中常见的摊销分析的例子，并且说明这个例子如何进行摊销分析
      \item 说明一下证明某个问题是NPC问题的两种方法，每种方法给出简要步骤
    \end{enumerate}

    \item 元素选择问题：输入为含有$n$个对象的集合$S$以及正整数$k$，要求返回$S$中第$k$小的元素。
    
    课程中提到了两种算法——基于分支算法与剪枝搜索的线性时间选择算法与随机选择算法。这两种算法的时间复杂度都是线性时间复杂度，请分别简化要说明这两种算法为什么能够实现线性时间选择（说明原因即可，不需要给出伪代码）。

    \item 考虑如下分布的数组$S$，其中元素先从小到大递增到$S[p]$，再从大到小递减。即$0 \sim P$范围单调递增，$p \sim \text{len}(S)-1$范围单调递减。要求设计一个时间复杂度为$O(\text{log}n)$的分治算法，找到这个位置$P$。请给出算法设计思路、伪代码以及时间复杂度的证明。

    \item 某工厂需要加工一批零件：$p_1, p_2, \cdots, p_n$，每个零件的加工时间为$a_1, a_2, \cdots, a_n$。该工厂一次只能加工一个零件，并且一个零件必须在一次内加工完成。假设这$n$个零件的加工顺序为：$p_{i1}, p_{i2}, \cdots p_{in}$，则每个零件的加工结束时间为：$e_{i1}, e_{i2}, \cdots, e_{in}$，请设计一个贪心算法，使得所有零件平均等待时间$(\sum_{k=1}^{n}{e_{ik}}) / n$最小化。请说明贪心选择性，给出伪代码，并分析复杂度。

	\item 考虑包含$n$个英文单词的文本排版问题。单词$W_k$的字符个数是$l_k$, $1 \leq k \leq n$。相邻两个单词之间有且仅有一个空格。这一段文本在显示屏上可能会分为多行，每一行容纳的字符数为$M$（包含空格）。请设计一个动态规划算法，要求最后一行前的所有行，最后一个单词的右侧的空白字符数最少（最后一行最右侧的空格不考虑）。计算空白个数的方法是：假设一行包含了第$i \sim j$个单词，那么该行最右侧留下的空白字符数就是：
	$$M - (j - i) - \sum_{k=i}^{j} l_k$$
	请：
	
	（1）给出算法的设计思路、优化子结构、重叠子问题 
	
	（2）给出伪代码，并分析复杂度。

	\item 说明课程中提到的素数测试算法、随机选择算法、随机快速排序算法分别属于哪种随机算法，并简要说明这几种随机算法的特点。
	
	\item 考虑如下通信问题：总共需要发送$n$个报文，每个报文的大小均不超过$C$，发送的方式是用容量为$C$的数据包打包。每个数据包可以容纳多个报文，只要这些报文的大小综合不超过$C$即可。请设计一个近似比为2的近似算法，完成上述通信任务。要求给出设计思路、伪代码、并证明该算法的近似比为2。

	\item 考虑基于分支限界剪枝搜索的TSP问题。在处理这个问题时，我们需要将TSP的邻接矩阵做如下处理：每行/每列都要减去当前行/列的最小元素，最终确保每行/每列都至少有一个0，最终构建得到代价矩阵。请证明，减去元素的综合可以作为TSP最短路径长度的一个代价下界。

	\item 考虑无向赋权图$G(V, E)$，其有一个结点子集$T \subseteq V$。我们需要找到一个具有最小权值总和的边集$E' \subseteq E$，使得$E'$能够连接$T$中的所有结点（即：$T$中的任意两个结点，都可以通过$E'$中的边相连）。请自选：模拟退火算法/遗传算法解决上述问题。要求：
	
	（1）不能只写出算法的框架，还要写出设计思路，以及(i) 对于模拟退火算法，写出新解的生成策略等； (ii) 对于遗传算法，写出染色体编码策略以及遗传操作（选择、交叉、变异等）。 此外，两种算法还要分析生成的解的合法性。
	
	（2）给出算法伪代码，并分析时间复杂度。

\end{enumerate}


\section{备注}

结合此前能够找到的两套试卷，不难发现考试题目大多都是课程算法案例的变形，甚至有些可以说是拙劣的变形，一眼就能看出是哪些已经讲过的典型问题，例如第4题装箱问题。所以需要好好听课+复习（尤其是要好好听课，否则复习起来很痛苦）。

考试难度并不大，但是题量很多。虽然只有9道题目，但每道题都需要写很多内容。

\textbf{关于回忆的题目的准确度：}

\begin{enumerate}
	\item 题目内容完整，无A/B卷划分
	\item 第4题所说的平均等待时间是：$(\sum_{k=1}^{n}{e_{ik}}) / n$，看起来有点奇怪，笔者认为应该是$(\sum_{k=1}^{n}{e_{ik}-a_{ik}}) / n$才对，但是试卷上确实写的是$(\sum_{k=1}^{n}{e_{ik}}) / n$。
	\item 第5题和第8题题目内容大致如此，可能有些细节记不清了，但是题意应该是类似的。
	\item 除此之外，其余题目应该与试卷原题保持一致，没有缺题。
\end{enumerate}

\textbf{个人评价：}

相比此前能够找到的两套试卷，今天考试时的试卷内容明显更侧重理解，而淡化了计算。明显的例子就是树搜索这一章，22年考察的是基于A*算法解决单源最短路径问题，要求画出搜索树的图即可。但是此份试卷考察的是分支限界算法求解TSP问题中的代价下界的理解。个人认为难度算是增加了，因为需要完全理解好这些算法才能从容应对（所以希望后来的学弟学妹们认真听课，虽然只是一门交叉课），当然回答答案也更加开放了。
\end{document}
